\section{Bring the formulas together}
\label{sec:comparison}
The three methods share query embedding, preparation, routing, filtering, and the optional
final reranker. Their difference is the work used to obtain candidates from the selected pool.
This is the right boundary for comparing the index algorithms without hiding common costs.

\Needspace{12\baselineskip}
\subsection{Latency, one method at a time}
Using the earlier definitions,
\begin{align}
T_A={}&\Tcommon+P c_{\rm range}
 +\Nelig Gc_{\rm lut}+H_s(\Nelig,R)c_{\rm heap}
 +T_{\rm rank}(R)+\Tfinish(R'),\label{eq:final-A}\\[3pt]
T_B={}&\Tcommon+T_{\rm order}+T_{\rm roots}
 +\sum_j s_jW_jc_{\rm split}+\sum_j\ell_jW_jc_{\rm leafword}\notag\\
 &+Vc_{\rm node}+FGc_{\rm lut,leaf}
 +H_s(F,R)c_{\rm heap}+T_{\rm rank}(R)+\Tfinish(R'),\label{eq:final-B}\\[3pt]
T_C={}&\Tcommon+T_{\rm keyprep}
 +\sum_j\left(E_jc_{\rm enum}+C_j^{\rm copy}c_{\rm copy}
              +H_jc_{\rm lookup}+F_jc_{\rm post}\right)\notag\\
 &+\sum_j F_jG_{C,j}c_{\rm lut,gather}+H_s(F,R)c_{\rm heap}
 +T_{\rm rank}(R)+\Tfinish(R').\label{eq:final-C}
\end{align}
$G_{C,j}$ is the candidate scoring work actually needed after reusing the key contribution; it can be zero for a full key with its exact score already known.

All cluster sums in query work run over $j\in\mathcal P(q)$. If C uses a merged stream,
add its explicitly counted stream-selection cost. Different execution choices can share
preparation or reuse buffers; record those choices rather than charging duplicated work.

\Needspace{12\baselineskip}
\subsection{An index speedup is not the same as an end-to-end speedup}
If two configurations have identical embedding and final-stage costs, write their shared time
as $T_0$. Replacing index time $t_A$ by $t_B$ gives speedup
\[
\frac{T_0+t_A}{T_0+t_B},
\]
not $t_A/t_B$. Even removing index time entirely cannot exceed $1+t_A/T_0$ when $T_0>0$.
If candidate counts or final-vector fetches differ, those costs are no longer shared and must
be evaluated separately. This is why cached-query index timings and complete text-query timings
should both be reported rather than substituted for one another.

\Needspace{12\baselineskip}
\subsection{Persistent storage}
For the reference layouts selected in this report,
\begin{align}
M_{A,\rm index}&=\Mbase,\label{eq:final-storageA}\\
M_{B,\rm index}&=\Mbase+d b_w\sum_{j=1}^{C}\ceil{n_j/w},\label{eq:final-storageB}\\
M_{C,\rm index}&=\Mbase+N b_{\rm ref}
 +\sum_{j=1}^{C}\ceil{U_j/\alpha}\,[b_k(h_j)+2b_o+b_{\rm slot}].
\label{eq:final-storageC}
\end{align}
Storage sums run over \emph{all} clusters, not only probed clusters. Increasing probes usually
changes query work and scratch, not the already built index size. For literal prefix relaxation, substitute the trie directory and query-work formula from Equation~\eqref{eq:trie-cost}.

A different C layout can
remove the additional posting permutation; it cannot remove the information needed to map
results back to documents.

\Needspace{12\baselineskip}
\subsection{Peak query RAM and the serving budget}
\fig{memory}{Persistent index data, runtime reserve and simultaneous scratch are separate terms. Sequential stage buffers need a lifetime analysis.}{.85}

The search-stage scratch terms are
\begin{align}
Q_{A,\rm search}&=\Qbase+Q_{\rm scan\ workspace},\\
Q_{B,\rm search}&=\Qbase+Q_{\rm order}
 +\max_t[\gamma_B Q_{B,\rm frontier}(t)+Q_{\rm current}(t)+Q_{\rm children}(t)],\\
Q_{C,\rm search}&=\Qbase+
 \max_t\left[\sum_j\gamma_C Q_j(t)b_{\rm enum,j}
 +Q_{\rm stream\ merge}(t)+Q_{\rm posting\ buffer}(t)+Q_{\rm gather}(t)\right].
\end{align}
Small iterator state is included in workspace terms. Embedding and final reranking can have
larger scratch than search, so use Equation~\eqref{eq:peak} to obtain the full query peak.
With resident index size $M_{m,\rm resident}$ for method $m$, the RAM constraint is
\begin{equation}
\boxed{M_{m,\rm resident}+M_{\rm model}+M_{\rm runtime}
       +\max_t\sum_{u\in\mathcal A(t)}Q_{m,u}(t)\leq M_{\rm RAM}.}
\label{eq:final-ram}
\end{equation}
$M_{\rm model}$ is the shared resident embedding-model memory if encoding runs locally; it is zero in this memory domain when an external service performs encoding. It is not multiplied by the number of queries.

For a conservative bound, replace the sum by $JQ_{m,\rm peak,max}$. A loose upper bound that
exceeds RAM does not prove that a tighter implementation cannot fit. An observed peak below
the limit does not prove that every possible query fits. These are different claims.

\Needspace{12\baselineskip}
\subsection{A compact comparison table}
\small
\begin{longtable}{@{}P{.16\linewidth}P{.25\linewidth}P{.25\linewidth}P{.25\linewidth}@{}}
\toprule & A: full scan & B: bitplanes & C: key enumeration\\\midrule\endfirsthead
\toprule & A: full scan & B: bitplanes & C: key enumeration\\\midrule\endhead
Additional index & none beyond common layout & $d b_w\sum_jW_j$ & postings and key directory\\
Principal work & every eligible code & split masks, leaf masks, full scores, frontier & empty and occupied keys, enumeration, postings, full scores\\
Variable counts & $\Nelig$ & $s_j,\ell_j,f_j,V$ & $H_j,F_j,E_j,Q_j^{\max}$\\
Largest scratch concern & selection or batch score buffers & many full-width pending masks & many pending flip sets or candidate gather buffers\\
Local exactness & scan all eligible rows & valid bounds, sufficient traversal, correct ties & full-key weighted order or a valid full-score bound\\
Approximation control & routing probes; candidate/rerank boundary & node budget, visit policy, earlier routing & prefix width, candidate or key budget, earlier routing\\
What must be measured & kernel cost, routing quality & skipped scores versus bitmap/queue cost and recall & empty-key volume versus useful candidates, timing and recall\\\bottomrule
\end{longtable}
\normalsize

\Needspace{12\baselineskip}
\subsection{The variables most likely to move the tradeoff}
\textbf{Documents per physical cluster, $n_j$.} A grows roughly linearly with eligible row
count. In B, a larger cluster also widens every active bitmap. In C, more documents per key can
reduce the number of empty probes needed, while increasing posting and scoring work.

\textbf{Binary dimension, $d$.} Longer codes increase score terms and stored bits. B may gain
more from avoiding full scores, but also stores more planes and can traverse more coordinates.
For full-key C, the number of possible keys grows exponentially. For short-key C, $h$ and $d$
are separate parameters.

\textbf{Probes, $P$.} More clusters can recover missed neighbors. They also increase eligible
rows, root masks, active key streams and candidate competition. Cluster sizes and the query's
routing behavior matter, not just the number $P$.

\textbf{Leaf size and node budget.} A small leaf threshold can trade scoring for more mask work.
A finite node budget caps visits but does not specify how many useful leaves are reached.
The best choice depends on the query and data, not a general monotonic recall formula.

\textbf{Lookup width and candidate target.} A short key has denser postings and a smaller key
space, but leaves more score information out. A longer key can improve candidate discrimination
and simultaneously make empty-key enumeration costly. A strong prefix can improve recall;
that does not remove the enumeration cost.

\textbf{Kernel and memory system.} SIMD, cache residency, batch size, alignment, layout and
thread count change the effective coefficients. Keep them fixed for an algorithm comparison,
or report their changes as a separate experiment.

\Needspace{12\baselineskip}
\subsection{Choosing configurations at fixed N and RAM}
Let $\theta_m$ contain the chosen method's tunable parameters and let $r_0$ be a recall target.
Choose the latency objective explicitly, for example median or p95 across queries. A valid
optimization statement is
\begin{equation}
\theta_m^*\in\arg\min_{\theta_m}\operatorname{Quantile}_{p}[T_m(q;\theta_m)]
\end{equation}
subject to the RAM condition in Equation~\eqref{eq:final-ram} and
\begin{equation}
\E_q[\operatorname{Recall@}K(q;\theta_m)]\geq r_0.
\label{eq:optimization}
\end{equation}
The expectation specifies an average-recall target. A requirement that almost every query
reach $r_0$ is stronger and needs a different probability constraint.

\par\noindent\begin{minipage}{\linewidth}
\begin{lstlisting}[caption={A defensible finite-grid experiment.}]
ChooseConfiguration(corpus, queries, fixed_reference, RAM_limit, targets):
    build or cache every configuration's required index
    compute exact reference results once for each query
    on tuning queries:
        for every configuration in the declared grid:
            measure latency, quality, work counts and peak RAM
            keep failures and empty returns in the record
    choose qualifying configurations for each target
    on separate test queries:
        evaluate those choices without retuning
    report lowest tested latency only where the target and RAM checks pass
    distinguish measured results from untested model predictions
\end{lstlisting}
\end{minipage}\par

A mathematical model can reduce the grid and explain its results. It cannot manufacture the
unknown query/document distribution or the true hardware coefficients. A ``best estimated
point'' from a model is a proposed experiment, not a verified engineering choice.

If $d$ is also varied, keep the quality reference fixed across choices, for example the same
long-vector top $K$ or the same relevance labels. Recomputing a different binary ground truth
for every $d$ makes a tiny representation look perfectly accurate against its own changed task.

\begin{resultbox}{Final interpretation}
A pays to score the selected pool. B pays to discover which scores can be avoided, and dense
bitmaps can make that discovery expensive. C pays to discover which keys contain useful
candidates, including empty keys and queue work. Each can be exact under a suitable stopping
rule. Which is fastest at a recall target is determined by the data, the representation,
the visited work, the memory layout and the hardware---not by one operation count or one
assumed normal distribution.
\end{resultbox}
